% Cuerpo de la tabla comparativa (se usa en la postulación y en tabla_comparacion.tex para el PNG)
\renewcommand{\arraystretch}{1.35}
\begin{tabularx}{\linewidth}{>{\raggedright\arraybackslash}p{4.9cm}>{\raggedright\arraybackslash}X>{\raggedright\arraybackslash}X}
\toprule
 & \textbf{Protocolo de Lo Barnechea} & \textbf{SAM-Cordillera} \\
\midrule
\multicolumn{3}{l}{\cellcolor{fondo}\textit{Qué entrega}} \\
Estados de la ruta & Tres alertas: preventiva, operativa y crítica & Seis estados: separa el cierre por despeje, por tormenta y por peligro \\
Escala espacial & Toda la ruta, con tres bandas fijas de isoterma & Cada tramo, con su cota, su exposición a avalanchas y su maquinaria; informa qué tramo determina el estado \\
Uso del pronóstico & Insumo de las reuniones D-7 y D-3, interpretado por los operadores & Entra a la misma regla que la observación y entrega el estado esperado por bloque de 6 horas \\
Peligro de avalancha y remoción & Evaluación técnica, sin umbral & Índice propio de 0 a 5, calculado aparte de la tormenta, con sus niveles de restricción, cierre parcial y cierre \\
Reapertura & Sin gatillantes explícitos & Se reabre cuando el peligro baja de su umbral y la nieve por despejar baja de 50~cm \\
Hora estimada de reapertura & No & Sí, cada 6 horas \\
Maquinaria & Estimación por escenario, fuera del cuadro de alertas & Equipos y operadores necesarios para reabrir en un plazo dado, cada 6 horas \\
Rutas donde se aplica & G-21 y sus accesos & Cualquier ruta con trazado, modelo de elevación y meteorología \\
\bottomrule
\end{tabularx}

\vspace{10pt}
\begin{tabularx}{\linewidth}{>{\raggedright\arraybackslash}p{4.3cm}>{\raggedright\arraybackslash}p{3.3cm}>{\raggedright\arraybackslash}X>{\raggedright\arraybackslash}X}
\toprule
 & \textbf{Decisión real}\newline{\footnotesize en la G-21, el protocolo aplicado por sus operadores} & \textbf{Umbrales escritos del protocolo}\newline{\footnotesize solo nieve, lluvia e isoterma} & \textbf{SAM-Cordillera} \\
\midrule
\multicolumn{4}{l}{\cellcolor{fondo}\textit{Cómo le fue en julio y agosto de 2026}} \\
G-21: períodos cerrados & \textbf{19} (101 horas) & justifican \textbf{0 de 19} & justifica \textbf{13 de 19} \\
60-CH: períodos cerrados & \textbf{176} (1.019 horas) & justifican \textbf{1 de 176} & justifica \textbf{72 de 176} \\
60-CH: períodos cerrados en que la regla indica operar con restricción o cierre parcial & -- & 142 de 176, sin evaluar peligro ni despeje & \textbf{77 de 176} \\
\bottomrule
\end{tabularx}

\vspace{4pt}
{\raggedright\footnotesize\color{txt2} Cada período es un bloque de 6 horas; julio y agosto suman 248 por ruta. La decisión real cuenta los períodos con al menos una hora de cierre por meteorología; en la 60-CH no se cuentan los 36 períodos cerrados por el horario del complejo fronterizo. En la G-21, esa decisión es la del protocolo de Lo Barnechea aplicado por sus operadores, que suman a los umbrales escritos el pronóstico y la evaluación técnica. Los umbrales escritos se revisan hora a hora; en la 60-CH ese protocolo no rige y se aplican solo como referencia.\par}
